% !TEX root = ../analisi-dei-capitolati.tex
% Capitolo 3 - Capitolato scelto: C2 - NEXUM Concierge

\section{Capitolato scelto: C2 -- NEXUM Concierge}\label{sec:c2}

\begin{scheda}
  \voce{Proponente}{Eggon S.r.l., Padova}
  \voce{Referente}{Gianpaolo Ferrarin}
  \voce{Dominio}{assistente conversazionale per i servizi HR di una
    piattaforma esistente}
\end{scheda}

\subsection{Obiettivo}
Eggon sviluppa \textbf{NEXUM}, una piattaforma per la comunicazione
aziendale e i servizi HR composta da un \glo{backoffice} amministrativo e
da una \glo{PWA} per il personale. Il progetto chiede di introdurre nella
PWA un \textbf{Concierge}: un assistente conversazionale con cui il
personale trova informazioni, capisce le procedure e usa i servizi
aziendali in linguaggio naturale, mentre dal backoffice l'azienda gestisce
le fonti e le capacità dell'assistente.

Per restare affrontabile, l'\glo{MVP} è limitato a \textbf{ferie e
permessi}, ma l'architettura deve permettere di aggiungere in futuro
altre fonti e abilità senza riprogettare il sistema. Il capitolato
prosegue la collaborazione avviata con l'Università nel 2025/2026, ma il
nuovo progetto ne è indipendente.

\subsection{Caratteristiche funzionali}
Il cuore del capitolato è riconoscere quattro tipi di richiesta, anche
all'interno della stessa conversazione (Tabella~\ref{tab:c2-capacita}).

\begin{table}[H]
  \centering
  \caption{Le quattro capacità del Concierge\label{tab:c2-capacita}}
  \begin{tabularx}{\linewidth}{@{}>{\raggedright\arraybackslash}p{3.3cm}
      >{\raggedright\arraybackslash}X >{\raggedright\arraybackslash}X@{}}
    \toprule
    \intestazione{Capacità} & \intestazione{Esempio} &
    \intestazione{Fonte o strumento} \\
    \midrule
    Conoscenza aziendale & «Come funzionano le ferie in azienda?» &
      Regolamenti caricati dal backoffice \\
    Conoscenza del prodotto & «Dove posso segnare le ferie?» & Manuali di
      NEXUM \\
    Dati personali & «Quante ferie ho fatto quest'anno?» & Abilità di
      lettura dei dati \\
    Operazioni & «Registra le mie ferie dal 10 al 15 ottobre» & Abilità di
      registrazione della richiesta \\
    \bottomrule
  \end{tabularx}
\end{table}

Attorno a queste capacità il BRD richiede:
\begin{itemize}
  \item una \textbf{chat testuale} che mantiene il contesto e chiede
        chiarimenti quando la richiesta è ambigua (per esempio una data
        senza anno);
  \item un \textbf{backoffice} per caricare, sostituire e ritirare le fonti
        (inizialmente PDF), organizzate in \glo{bucket} con destinatari
        autorizzati; le risposte devono citare le fonti e il sistema deve
        astenersi quando le informazioni non bastano;
  \item \textbf{dati personali} letti solo tramite i servizi di NEXUM: il
        modello non deve inventare dati né fare conteggi al posto dei
        servizi;
  \item un \textbf{riepilogo completo e una conferma esplicita} prima di
        ogni invio, distinguendo tra richiesta registrata e approvata;
  \item \textbf{autorizzazioni, gestione degli errori e tracciabilità}:
        due profili (personale e amministrazione), separazione tra aziende,
        registrazione delle fonti usate e degli esiti delle operazioni.
\end{itemize}

\subsection{Requisiti e criteri di accettazione}
Il BRD elenca sedici requisiti di business, che saranno la base
dell'Analisi dei Requisiti. Non distingue requisiti obbligatori e
opzionali: il nucleo obbligatorio è l'intero percorso sul dominio scelto,
mentre grafi di conoscenza, distillazione avanzata, altri formati e altri
domini sono possibilità ulteriori. Per l'accettazione l'MVP deve
dimostrare tutte e quattro le capacità e un percorso unitario dalla
domanda informativa all'invio della richiesta, verificando inoltre
chiarimenti, citazioni e astensione, aggiornamento dopo il ritiro di una
fonte, autorizzazioni e gestione degli errori.

Il lavoro è diviso in tre fasi: un \glo{PoC} con le interfacce concordate
(alla RTB), poi il percorso completo e la verifica dell'MVP (alla PB).
L'MVP può usare servizi NEXUM di prova o simulatori conformi al contratto
di integrazione definito da Eggon, purché il percorso sia eseguibile da
capo a fondo.

\subsection{Considerazioni tecnologiche}
Il capitolato non impone tecnologie. Nell'incontro Eggon ha indicato come
riferimento non vincolante il proprio stack: Ruby on Rails per il backend,
Angular per backoffice e PWA, RubyLLM come interfaccia unica verso i
provider di modelli (per esempio Amazon Bedrock, Google Vertex,
OpenRouter), infrastruttura AWS. Chiede però un approccio
\textbf{agnostico rispetto al provider} e una scelta motivata da
esperimenti:
\begin{itemize}
  \item per la \textbf{conoscenza documentale}, confrontare il documento
        inserito nel contesto della conversazione con soluzioni basate su
        \glo{embedding}, ricerca vettoriale ed eventualmente grafi
        (\glo{RAG});
  \item per le \textbf{operazioni}, esporre le abilità tramite \glo{MCP},
        con un registro che ne definisce ingressi e uscite e un
        classificatore che riconosce l'intenzione dell'utente (posizione
        ferie, richiesta, conferma, annullamento, spostamento).
\end{itemize}
Il gruppo partirà da servizi simulati per poi passare ai modelli reali su
AWS, controllando il consumo di token.

\subsection{Incontro con il proponente}
Dall'incontro del 26 ottobre 2026 (VE~2026-10-26) sono emersi questi
chiarimenti:
\begin{itemize}
  \item \textbf{perimetro}: MVP solo su ferie e permessi, con il
        regolamento aziendale in PDF come primo documento e solo lingua
        italiana;
  \item \textbf{ricerca applicata}: il valore atteso sta nel confronto tra
        approcci e nella qualità delle decisioni, non solo nel PoC;
  \item \textbf{risorse}: accesso ad AWS e ai modelli con limiti di
        budget, i due repository dei gruppi precedenti come riferimento,
        esempi e materiali di prova costruiti insieme;
  \item \textbf{dati}: solo dati fittizi in sperimentazione;
  \item \textbf{gruppo}: sei persone invece delle sette del BRD non sono un
        problema;
  \item \textbf{collaborazione}: incontro settimanale di circa un'ora in un
        giorno fisso, Telegram per le comunicazioni rapide, Meet per gli
        incontri; Eggon chiede un confronto proattivo sulle scelte, anche
        nei periodi di esami;
  \item \textbf{strumenti di generazione del codice}: ammessi, purché il
        gruppo sappia spiegare il codice e ne curi la sicurezza (OWASP,
        Cyber Resilience Act); Eggon è disponibile a mostrare come lavora
        con Claude Code.
\end{itemize}

\subsection{Rischi principali}
\begin{itemize}
  \item \textbf{Ampliamento del perimetro}: mantenere ferie e permessi come
        unico dominio dell'MVP.
  \item \textbf{Risposte non supportate dalle fonti}: citazioni
        verificabili, astensione e un insieme di domande di prova per
        misurare la qualità.
  \item \textbf{Invii errati o duplicati}: riepilogo, conferma esplicita e
        verifica dell'esito prima di ripetere un'operazione.
  \item \textbf{Dipendenza dalle integrazioni}: concordare presto il
        contratto con Eggon e sviluppare su simulatori conformi.
  \item \textbf{Tecnologie nuove}: studio mirato nelle prime settimane e
        incontri di approfondimento con il proponente.
  \item \textbf{Continuità della comunicazione}, indicata da Eggon come
        rischio principale: incontro settimanale fisso anche nei periodi di
        esami.
  \item \textbf{Costi dei modelli}: servizi simulati all'inizio e
        monitoraggio del consumo di token entro il budget concordato.
\end{itemize}

\subsection{Valutazione del gruppo}

\begin{fortedeboli}
  \fd{\textbf{Perimetro circoscritto}: un solo dominio e vincoli espliciti
      su ciò che non è richiesto.}
     {\textbf{Scelte da sperimentare}: stack e strategia documentale vanno
      decisi e motivati con esperimenti.}
  \fd{\textbf{Criteri di accettazione verificabili}: ottima base per
      requisiti e test.}
     {\textbf{Qualità delle risposte difficile da misurare}: servono
      metriche e domande di prova costruite dal gruppo.}
  \fd{\textbf{Tecnologie attuali}: RAG, MCP e uso di strumenti da parte
      di un \glo{LLM}.}
     {\textbf{Tecnologie nuove per il gruppo}: serve un investimento
      iniziale di studio.}
  \fd{\textbf{Prodotto reale e proponente presente}: infrastruttura
      fornita, incontri settimanali, materiale dei gruppi precedenti.}
     {\textbf{Dipendenze}: contratto di integrazione da Eggon e budget
      per i modelli da rispettare.}
\end{fortedeboli}

Il gruppo ha scelto C2 per l'equilibrio tra interesse, valore formativo e
sostenibilità: un tema attuale in un prodotto reale, con un perimetro
limitato e criteri verificabili, senza hardware né tecnologie imposte.
L'incontro con Eggon ha confermato la scelta: il proponente fornisce le
risorse necessarie e cerca proprio una collaborazione basata sulla
sperimentazione e sul confronto delle scelte.
